\documentclass[11pt,letterpaper]{article}
\usepackage[technical-reference]{cornell-notes}
\usepackage{needspace}

\title{Clause 9: Declarations - Cornell Notes}
\author{}
\date{}
\setDocTitle{Clause 9: Declarations}
\setDocSubtitle{C++ 2024 Cornell Notes}
\setDocOwner{C++ 2024 Cornell Notes}
\setDocAuthor{}
\setDocDate{}
\setCornellCollection{C++ 2024 Cornell Notes}
\setCornellUnitType{Clause}
\setCornellUnitNumber{9}
\setCornellUnitTitle{Declarations}
\setCornellCompanionLabel{C++ Technical Reference}
\hypersetup{pdftitle={Clause 9: Declarations - Cornell Notes},pdfsubject={C++ technical study notes},pdfauthor={}}

\begin{document}
\maketitle
\begin{CornellReferenceOrientationBox}
\textbf{Purpose.} Defines specifiers, declarators, initialization, function definitions, structured bindings, enumerations, namespaces, using declarations, linkage specifications, inline assembly declarations, and standard attributes.
\par\textbf{Role.} Creates and describes entities, their types, initialization, and program-visible properties.
\par\textbf{Principal dependencies.} Uses scope, lookup, types, expressions, and statements; supplies entities to modules, classes, overloading, and templates.
\par\textbf{Primary audience.} programmers, compilers, reflection-like tools, and API reviewers.
\par\textbf{Major subject groups.} Preamble; Specifiers; Declarators; Initializers; Function definitions; Structured binding declarations; Enumerations; Namespaces; The using declaration; The asm declaration; Linkage specifications; Attributes.
\end{CornellReferenceOrientationBox}
\section{Learning Objectives}
\begin{itemize}[label=\textcolor{CornellReferenceTeal}{$\blacktriangleright$}]
\item Define the governing ideas of Preamble and state the consequences for a program or implementation.
\item Identify the governing ideas of Specifiers and state the consequences for a program or implementation.
\item Distinguish the governing ideas of Declarators and state the consequences for a program or implementation.
\item Explain the governing ideas of Initializers and state the consequences for a program or implementation.
\item Trace the governing ideas of Function definitions and state the consequences for a program or implementation.
\item Apply the governing ideas of Structured binding declarations and state the consequences for a program or implementation.
\item Analyze the governing ideas of Enumerations and state the consequences for a program or implementation.
\item Evaluate the governing ideas of Namespaces and state the consequences for a program or implementation.
\item Diagnose the governing ideas of The using declaration and state the consequences for a program or implementation.
\item Compare the governing ideas of The asm declaration and state the consequences for a program or implementation.
\item Verify the governing ideas of Linkage specifications and state the consequences for a program or implementation.
\item Relate the governing ideas of Attributes and state the consequences for a program or implementation.
\item Classify the governing ideas of General and state the consequences for a program or implementation.
\item Justify the governing ideas of Storage class specifiers and state the consequences for a program or implementation.
\item Audit the governing ideas of Function specifiers and state the consequences for a program or implementation.
\item Summarize the governing ideas of The typedef specifier and state the consequences for a program or implementation.
\end{itemize}
\section{Normative Structure Map}
\small\begin{longtable}{@{}>{\RaggedRight\arraybackslash}p{0.13\textwidth}>{\RaggedRight\arraybackslash}p{0.27\textwidth}>{\RaggedRight\arraybackslash}p{0.19\textwidth}>{\RaggedRight\arraybackslash}p{0.31\textwidth}@{}}
\toprule\textbf{Subclause} & \textbf{Subject} & \textbf{Rule category} & \textbf{Key dependencies / entities}\\\midrule\endfirsthead
\toprule\textbf{Subclause} & \textbf{Subject} & \textbf{Rule category} & \textbf{Key dependencies / entities}\\\midrule\endhead
\bottomrule\endfoot
9.1 & Preamble & core-language semantics & Uses scope, lookup, types, expressions, and statements; supplies entities to modules, classes, overloading, and templates.\\
9.2 & Specifiers & core-language semantics & General, Storage class specifiers, Function specifiers, The typedef specifier, and related subclauses\\
9.3 & Declarators & syntax / lexical rules & General, Type names, Ambiguity resolution, Meaning of declarators\\
9.4 & Initializers & core-language semantics & General, Aggregates, Character arrays, References, and related subclauses\\
9.5 & Function definitions & operation semantics & General, Explicitly-defaulted functions, Deleted definitions, Coroutine definitions\\
9.6 & Structured binding declarations & core-language semantics & Uses scope, lookup, types, expressions, and statements; supplies entities to modules, classes, overloading, and templates.\\
9.7 & Enumerations & core-language semantics & Enumeration declarations, The using enum declaration\\
9.8 & Namespaces & core-language semantics & General, Namespace definition, Namespace alias, Using namespace directive\\
9.9 & The using declaration & core-language semantics & Uses scope, lookup, types, expressions, and statements; supplies entities to modules, classes, overloading, and templates.\\
9.10 & The asm declaration & core-language semantics & Uses scope, lookup, types, expressions, and statements; supplies entities to modules, classes, overloading, and templates.\\
9.11 & Linkage specifications & core-language semantics & Uses scope, lookup, types, expressions, and statements; supplies entities to modules, classes, overloading, and templates.\\
9.12 & Attributes & core-language semantics & Attribute syntax and semantics, Alignment specifier, Assumption attribute, Carries dependency attribute, and related subclauses\\
\end{longtable}\normalsize
\begin{CornellReferenceNormativeBox}A heading maps the clause; it does not replace the detailed conditions. For each operation, read requirements, effects, result, exceptions, complexity, remarks, and notes with their stated normative force.\end{CornellReferenceNormativeBox}
\subsection*{Rule Classification Guide}
\small\begin{longtable}{@{}>{\RaggedRight\arraybackslash}p{0.25\textwidth}>{\RaggedRight\arraybackslash}p{0.67\textwidth}@{}}\toprule\textbf{Label or category} & \textbf{How to read it}\\\midrule\endfirsthead\toprule\textbf{Label or category} & \textbf{How to read it (continued)}\\\midrule\endhead\bottomrule\endfoot
Well-formed / ill-formed & Formation is governed by syntax and semantic rules; an ill-formed program does not become well-formed because one implementation accepts it.\\
Diagnosable rule & A violation requires at least one diagnostic unless the wording places the case outside the diagnosable set.\\
No diagnostic required & The program can be ill-formed while the implementation has no diagnostic obligation.\\
Undefined behavior & No execution requirements are imposed; translation success does not establish safety or portability.\\
Unspecified behavior & One of the permitted outcomes may be chosen without documentation.\\
Implementation-defined behavior & The implementation chooses and documents the behavior.\\
Conditionally supported & The implementation may omit support but documents unsupported constructs.\\
\end{longtable}\normalsize
\section{Cornell Cue-and-Notes}
\Needspace{8\baselineskip}
\subsection*{9.1 Preamble}
\begin{longtable}{@{}>{\RaggedRight\arraybackslash\columncolor{CornellReferenceCueGray}}p{0.28\textwidth}>{\RaggedRight\arraybackslash}p{0.64\textwidth}@{}}
\rowcolor{CornellReferenceNavy}\color{white}\textbf{Cue / recall prompt} & \color{white}\textbf{Notes}\\\endfirsthead
\rowcolor{CornellReferenceNavy}\color{white}\textbf{Cue / recall prompt} & \color{white}\textbf{Notes (continued)}\\\endhead
\arrayrulecolor{CornellReferenceLineGray}
\textbf{What rule applies to Preamble?}\\[-1mm]{\scriptsize\CornellClauseRef{9.1}} & Frames the terminology, applicability, and relationships needed to interpret Preamble; detailed rules remain in the following subclauses.\\\hline
\end{longtable}
\begin{CornellReferenceSummaryBox}Preamble supplies the core-language semantics portion of this clause. The key review move is to connect its local wording to the clause-wide model, then classify each condition by normative force before relying on it.\end{CornellReferenceSummaryBox}
\Needspace{8\baselineskip}
\subsection*{9.2 Specifiers}
\begin{longtable}{@{}>{\RaggedRight\arraybackslash\columncolor{CornellReferenceCueGray}}p{0.28\textwidth}>{\RaggedRight\arraybackslash}p{0.64\textwidth}@{}}
\rowcolor{CornellReferenceNavy}\color{white}\textbf{Cue / recall prompt} & \color{white}\textbf{Notes}\\\endfirsthead
\rowcolor{CornellReferenceNavy}\color{white}\textbf{Cue / recall prompt} & \color{white}\textbf{Notes (continued)}\\\endhead
\arrayrulecolor{CornellReferenceLineGray}
\textbf{What rule applies to General?}\\[-1mm]{\scriptsize\CornellClauseRef{9.2.1}} & Frames the terminology, applicability, and relationships needed to interpret General; detailed rules remain in the following subclauses.\\\hline
\textbf{What rule applies to Storage class specifiers?}\\[-1mm]{\scriptsize\CornellClauseRef{9.2.2}} & Defines the core-language rules for Storage class specifiers, including formation, semantic effect, interactions with type and lookup, and any ill-formed or implementation-dependent cases.\\\hline
\textbf{What rule applies to Function specifiers?}\\[-1mm]{\scriptsize\CornellClauseRef{9.2.3}} & Defines the core-language rules for Function specifiers, including formation, semantic effect, interactions with type and lookup, and any ill-formed or implementation-dependent cases.\\\hline
\textbf{What rule applies to The typedef specifier?}\\[-1mm]{\scriptsize\CornellClauseRef{9.2.4}} & Defines the core-language rules for The typedef specifier, including formation, semantic effect, interactions with type and lookup, and any ill-formed or implementation-dependent cases.\\\hline
\textbf{What rule applies to The friend specifier?}\\[-1mm]{\scriptsize\CornellClauseRef{9.2.5}} & Defines the core-language rules for The friend specifier, including formation, semantic effect, interactions with type and lookup, and any ill-formed or implementation-dependent cases.\\\hline
\textbf{What rule applies to The constexpr and consteval specifiers?}\\[-1mm]{\scriptsize\CornellClauseRef{9.2.6}} & Defines the core-language rules for The constexpr and consteval specifiers, including formation, semantic effect, interactions with type and lookup, and any ill-formed or implementation-dependent cases.\\\hline
\textbf{What rule applies to The constinit specifier?}\\[-1mm]{\scriptsize\CornellClauseRef{9.2.7}} & Defines the core-language rules for The constinit specifier, including formation, semantic effect, interactions with type and lookup, and any ill-formed or implementation-dependent cases.\\\hline
\textbf{What rule applies to The inline specifier?}\\[-1mm]{\scriptsize\CornellClauseRef{9.2.8}} & Defines the core-language rules for The inline specifier, including formation, semantic effect, interactions with type and lookup, and any ill-formed or implementation-dependent cases.\\\hline
\textbf{What rule applies to Type specifiers?}\\[-1mm]{\scriptsize\CornellClauseRef{9.2.9}} & Defines the core-language rules for Type specifiers, including formation, semantic effect, interactions with type and lookup, and any ill-formed or implementation-dependent cases.\\\hline
\end{longtable}
\begin{CornellReferenceSummaryBox}Specifiers supplies the core-language semantics portion of this clause. The key review move is to connect its local wording to the clause-wide model, then classify each condition by normative force before relying on it.\end{CornellReferenceSummaryBox}
\Needspace{5\baselineskip}\paragraph{Detailed coverage ledger.}
\begin{description}[style=nextline,leftmargin=2.8cm,font=\normalfont\bfseries\color{CornellReferenceTeal}]
\item[9.2.9.1] \textbf{General.} Frames the terminology, applicability, and relationships needed to interpret General; detailed rules remain in the following subclauses.
\item[9.2.9.2] \textbf{The cv-qualifiers.} Defines the core-language rules for The cv-qualifiers, including formation, semantic effect, interactions with type and lookup, and any ill-formed or implementation-dependent cases.
\item[9.2.9.3] \textbf{Simple type specifiers.} Defines the core-language rules for Simple type specifiers, including formation, semantic effect, interactions with type and lookup, and any ill-formed or implementation-dependent cases.
\item[9.2.9.4] \textbf{Elaborated type specifiers.} Defines the core-language rules for Elaborated type specifiers, including formation, semantic effect, interactions with type and lookup, and any ill-formed or implementation-dependent cases.
\item[9.2.9.5] \textbf{Decltype specifiers.} Defines the core-language rules for Decltype specifiers, including formation, semantic effect, interactions with type and lookup, and any ill-formed or implementation-dependent cases.
\item[9.2.9.6] \textbf{Placeholder type specifiers.} Defines the core-language rules for Placeholder type specifiers, including formation, semantic effect, interactions with type and lookup, and any ill-formed or implementation-dependent cases.
\item[9.2.9.6.1] \textbf{General.} Frames the terminology, applicability, and relationships needed to interpret General; detailed rules remain in the following subclauses.
\item[9.2.9.6.2] \textbf{Placeholder type deduction.} Defines the template context for Placeholder type deduction; track dependence, lookup time, substitution, constraints, instantiation, and specialization ordering separately.
\item[9.2.9.7] \textbf{Deduced class template specialization types.} Defines the template context for Deduced class template specialization types; track dependence, lookup time, substitution, constraints, instantiation, and specialization ordering separately.
\end{description}
\Needspace{8\baselineskip}
\subsection*{9.3 Declarators}
\begin{longtable}{@{}>{\RaggedRight\arraybackslash\columncolor{CornellReferenceCueGray}}p{0.28\textwidth}>{\RaggedRight\arraybackslash}p{0.64\textwidth}@{}}
\rowcolor{CornellReferenceNavy}\color{white}\textbf{Cue / recall prompt} & \color{white}\textbf{Notes}\\\endfirsthead
\rowcolor{CornellReferenceNavy}\color{white}\textbf{Cue / recall prompt} & \color{white}\textbf{Notes (continued)}\\\endhead
\arrayrulecolor{CornellReferenceLineGray}
\textbf{What rule applies to General?}\\[-1mm]{\scriptsize\CornellClauseRef{9.3.1}} & Frames the terminology, applicability, and relationships needed to interpret General; detailed rules remain in the following subclauses.\\\hline
\textbf{What rule applies to Type names?}\\[-1mm]{\scriptsize\CornellClauseRef{9.3.2}} & Defines the core-language rules for Type names, including formation, semantic effect, interactions with type and lookup, and any ill-formed or implementation-dependent cases.\\\hline
\textbf{What rule applies to Ambiguity resolution?}\\[-1mm]{\scriptsize\CornellClauseRef{9.3.3}} & Defines the core-language rules for Ambiguity resolution, including formation, semantic effect, interactions with type and lookup, and any ill-formed or implementation-dependent cases.\\\hline
\textbf{What rule applies to Meaning of declarators?}\\[-1mm]{\scriptsize\CornellClauseRef{9.3.4}} & Defines the core-language rules for Meaning of declarators, including formation, semantic effect, interactions with type and lookup, and any ill-formed or implementation-dependent cases.\\\hline
\end{longtable}
\begin{CornellReferenceSummaryBox}Declarators supplies the syntax / lexical rules portion of this clause. The key review move is to connect its local wording to the clause-wide model, then classify each condition by normative force before relying on it.\end{CornellReferenceSummaryBox}
\Needspace{5\baselineskip}\paragraph{Detailed coverage ledger.}
\begin{description}[style=nextline,leftmargin=2.8cm,font=\normalfont\bfseries\color{CornellReferenceTeal}]
\item[9.3.4.1] \textbf{General.} Frames the terminology, applicability, and relationships needed to interpret General; detailed rules remain in the following subclauses.
\item[9.3.4.2] \textbf{Pointers.} Defines the core-language rules for Pointers, including formation, semantic effect, interactions with type and lookup, and any ill-formed or implementation-dependent cases.
\item[9.3.4.3] \textbf{References.} Defines the core-language rules for References, including formation, semantic effect, interactions with type and lookup, and any ill-formed or implementation-dependent cases.
\item[9.3.4.4] \textbf{Pointers to members.} Defines the core-language rules for Pointers to members, including formation, semantic effect, interactions with type and lookup, and any ill-formed or implementation-dependent cases.
\item[9.3.4.5] \textbf{Arrays.} Defines the core-language rules for Arrays, including formation, semantic effect, interactions with type and lookup, and any ill-formed or implementation-dependent cases.
\item[9.3.4.6] \textbf{Functions.} Defines the core-language rules for Functions, including formation, semantic effect, interactions with type and lookup, and any ill-formed or implementation-dependent cases.
\item[9.3.4.7] \textbf{Default arguments.} Defines the core-language rules for Default arguments, including formation, semantic effect, interactions with type and lookup, and any ill-formed or implementation-dependent cases.
\end{description}
\Needspace{8\baselineskip}
\subsection*{9.4 Initializers}
\begin{longtable}{@{}>{\RaggedRight\arraybackslash\columncolor{CornellReferenceCueGray}}p{0.28\textwidth}>{\RaggedRight\arraybackslash}p{0.64\textwidth}@{}}
\rowcolor{CornellReferenceNavy}\color{white}\textbf{Cue / recall prompt} & \color{white}\textbf{Notes}\\\endfirsthead
\rowcolor{CornellReferenceNavy}\color{white}\textbf{Cue / recall prompt} & \color{white}\textbf{Notes (continued)}\\\endhead
\arrayrulecolor{CornellReferenceLineGray}
\textbf{What rule applies to General?}\\[-1mm]{\scriptsize\CornellClauseRef{9.4.1}} & Frames the terminology, applicability, and relationships needed to interpret General; detailed rules remain in the following subclauses.\\\hline
\textbf{What rule applies to Aggregates?}\\[-1mm]{\scriptsize\CornellClauseRef{9.4.2}} & Defines the core-language rules for Aggregates, including formation, semantic effect, interactions with type and lookup, and any ill-formed or implementation-dependent cases.\\\hline
\textbf{What rule applies to Character arrays?}\\[-1mm]{\scriptsize\CornellClauseRef{9.4.3}} & Specifies text-sequence behavior for Character arrays; establish character type, extent, ownership, null termination, encoding assumptions, and invalidation before use.\\\hline
\textbf{What rule applies to References?}\\[-1mm]{\scriptsize\CornellClauseRef{9.4.4}} & Defines the core-language rules for References, including formation, semantic effect, interactions with type and lookup, and any ill-formed or implementation-dependent cases.\\\hline
\textbf{What rule applies to List-initialization?}\\[-1mm]{\scriptsize\CornellClauseRef{9.4.5}} & Determines the initialization form used by List-initialization, the candidates or conversions considered, object-lifetime start, narrowing rules, and failure consequences.\\\hline
\end{longtable}
\begin{CornellReferenceSummaryBox}Initializers supplies the core-language semantics portion of this clause. The key review move is to connect its local wording to the clause-wide model, then classify each condition by normative force before relying on it.\end{CornellReferenceSummaryBox}
\Needspace{8\baselineskip}
\subsection*{9.5 Function definitions}
\begin{longtable}{@{}>{\RaggedRight\arraybackslash\columncolor{CornellReferenceCueGray}}p{0.28\textwidth}>{\RaggedRight\arraybackslash}p{0.64\textwidth}@{}}
\rowcolor{CornellReferenceNavy}\color{white}\textbf{Cue / recall prompt} & \color{white}\textbf{Notes}\\\endfirsthead
\rowcolor{CornellReferenceNavy}\color{white}\textbf{Cue / recall prompt} & \color{white}\textbf{Notes (continued)}\\\endhead
\arrayrulecolor{CornellReferenceLineGray}
\textbf{What rule applies to General?}\\[-1mm]{\scriptsize\CornellClauseRef{9.5.1}} & Frames the terminology, applicability, and relationships needed to interpret General; detailed rules remain in the following subclauses.\\\hline
\textbf{What rule applies to Explicitly-defaulted functions?}\\[-1mm]{\scriptsize\CornellClauseRef{9.5.2}} & Defines the core-language rules for Explicitly-defaulted functions, including formation, semantic effect, interactions with type and lookup, and any ill-formed or implementation-dependent cases.\\\hline
\textbf{What rule applies to Deleted definitions?}\\[-1mm]{\scriptsize\CornellClauseRef{9.5.3}} & Defines the core-language rules for Deleted definitions, including formation, semantic effect, interactions with type and lookup, and any ill-formed or implementation-dependent cases.\\\hline
\textbf{What rule applies to Coroutine definitions?}\\[-1mm]{\scriptsize\CornellClauseRef{9.5.4}} & Defines the core-language rules for Coroutine definitions, including formation, semantic effect, interactions with type and lookup, and any ill-formed or implementation-dependent cases.\\\hline
\end{longtable}
\begin{CornellReferenceSummaryBox}Function definitions supplies the operation semantics portion of this clause. The key review move is to connect its local wording to the clause-wide model, then classify each condition by normative force before relying on it.\end{CornellReferenceSummaryBox}
\Needspace{8\baselineskip}
\subsection*{9.6 Structured binding declarations}
\begin{longtable}{@{}>{\RaggedRight\arraybackslash\columncolor{CornellReferenceCueGray}}p{0.28\textwidth}>{\RaggedRight\arraybackslash}p{0.64\textwidth}@{}}
\rowcolor{CornellReferenceNavy}\color{white}\textbf{Cue / recall prompt} & \color{white}\textbf{Notes}\\\endfirsthead
\rowcolor{CornellReferenceNavy}\color{white}\textbf{Cue / recall prompt} & \color{white}\textbf{Notes (continued)}\\\endhead
\arrayrulecolor{CornellReferenceLineGray}
\textbf{What rule applies to Structured binding declarations?}\\[-1mm]{\scriptsize\CornellClauseRef{9.6}} & Defines the core-language rules for Structured binding declarations, including formation, semantic effect, interactions with type and lookup, and any ill-formed or implementation-dependent cases.\\\hline
\end{longtable}
\begin{CornellReferenceSummaryBox}Structured binding declarations supplies the core-language semantics portion of this clause. The key review move is to connect its local wording to the clause-wide model, then classify each condition by normative force before relying on it.\end{CornellReferenceSummaryBox}
\Needspace{8\baselineskip}
\subsection*{9.7 Enumerations}
\begin{longtable}{@{}>{\RaggedRight\arraybackslash\columncolor{CornellReferenceCueGray}}p{0.28\textwidth}>{\RaggedRight\arraybackslash}p{0.64\textwidth}@{}}
\rowcolor{CornellReferenceNavy}\color{white}\textbf{Cue / recall prompt} & \color{white}\textbf{Notes}\\\endfirsthead
\rowcolor{CornellReferenceNavy}\color{white}\textbf{Cue / recall prompt} & \color{white}\textbf{Notes (continued)}\\\endhead
\arrayrulecolor{CornellReferenceLineGray}
\textbf{What rule applies to Enumeration declarations?}\\[-1mm]{\scriptsize\CornellClauseRef{9.7.1}} & Defines the core-language rules for Enumeration declarations, including formation, semantic effect, interactions with type and lookup, and any ill-formed or implementation-dependent cases.\\\hline
\textbf{What rule applies to The using enum declaration?}\\[-1mm]{\scriptsize\CornellClauseRef{9.7.2}} & Defines the core-language rules for The using enum declaration, including formation, semantic effect, interactions with type and lookup, and any ill-formed or implementation-dependent cases.\\\hline
\end{longtable}
\begin{CornellReferenceSummaryBox}Enumerations supplies the core-language semantics portion of this clause. The key review move is to connect its local wording to the clause-wide model, then classify each condition by normative force before relying on it.\end{CornellReferenceSummaryBox}
\Needspace{8\baselineskip}
\subsection*{9.8 Namespaces}
\begin{longtable}{@{}>{\RaggedRight\arraybackslash\columncolor{CornellReferenceCueGray}}p{0.28\textwidth}>{\RaggedRight\arraybackslash}p{0.64\textwidth}@{}}
\rowcolor{CornellReferenceNavy}\color{white}\textbf{Cue / recall prompt} & \color{white}\textbf{Notes}\\\endfirsthead
\rowcolor{CornellReferenceNavy}\color{white}\textbf{Cue / recall prompt} & \color{white}\textbf{Notes (continued)}\\\endhead
\arrayrulecolor{CornellReferenceLineGray}
\textbf{What rule applies to General?}\\[-1mm]{\scriptsize\CornellClauseRef{9.8.1}} & Frames the terminology, applicability, and relationships needed to interpret General; detailed rules remain in the following subclauses.\\\hline
\textbf{What rule applies to Namespace definition?}\\[-1mm]{\scriptsize\CornellClauseRef{9.8.2}} & Defines the core-language rules for Namespace definition, including formation, semantic effect, interactions with type and lookup, and any ill-formed or implementation-dependent cases.\\\hline
\textbf{What rule applies to Namespace alias?}\\[-1mm]{\scriptsize\CornellClauseRef{9.8.3}} & Defines the core-language rules for Namespace alias, including formation, semantic effect, interactions with type and lookup, and any ill-formed or implementation-dependent cases.\\\hline
\textbf{What rule applies to Using namespace directive?}\\[-1mm]{\scriptsize\CornellClauseRef{9.8.4}} & Defines the core-language rules for Using namespace directive, including formation, semantic effect, interactions with type and lookup, and any ill-formed or implementation-dependent cases.\\\hline
\end{longtable}
\begin{CornellReferenceSummaryBox}Namespaces supplies the core-language semantics portion of this clause. The key review move is to connect its local wording to the clause-wide model, then classify each condition by normative force before relying on it.\end{CornellReferenceSummaryBox}
\Needspace{5\baselineskip}\paragraph{Detailed coverage ledger.}
\begin{description}[style=nextline,leftmargin=2.8cm,font=\normalfont\bfseries\color{CornellReferenceTeal}]
\item[9.8.2.1] \textbf{General.} Frames the terminology, applicability, and relationships needed to interpret General; detailed rules remain in the following subclauses.
\item[9.8.2.2] \textbf{Unnamed namespaces.} Defines the core-language rules for Unnamed namespaces, including formation, semantic effect, interactions with type and lookup, and any ill-formed or implementation-dependent cases.
\end{description}
\Needspace{8\baselineskip}
\subsection*{9.9 The using declaration}
\begin{longtable}{@{}>{\RaggedRight\arraybackslash\columncolor{CornellReferenceCueGray}}p{0.28\textwidth}>{\RaggedRight\arraybackslash}p{0.64\textwidth}@{}}
\rowcolor{CornellReferenceNavy}\color{white}\textbf{Cue / recall prompt} & \color{white}\textbf{Notes}\\\endfirsthead
\rowcolor{CornellReferenceNavy}\color{white}\textbf{Cue / recall prompt} & \color{white}\textbf{Notes (continued)}\\\endhead
\arrayrulecolor{CornellReferenceLineGray}
\textbf{What rule applies to The using declaration?}\\[-1mm]{\scriptsize\CornellClauseRef{9.9}} & Defines the core-language rules for The using declaration, including formation, semantic effect, interactions with type and lookup, and any ill-formed or implementation-dependent cases.\\\hline
\end{longtable}
\begin{CornellReferenceSummaryBox}The using declaration supplies the core-language semantics portion of this clause. The key review move is to connect its local wording to the clause-wide model, then classify each condition by normative force before relying on it.\end{CornellReferenceSummaryBox}
\Needspace{8\baselineskip}
\subsection*{9.10 The asm declaration}
\begin{longtable}{@{}>{\RaggedRight\arraybackslash\columncolor{CornellReferenceCueGray}}p{0.28\textwidth}>{\RaggedRight\arraybackslash}p{0.64\textwidth}@{}}
\rowcolor{CornellReferenceNavy}\color{white}\textbf{Cue / recall prompt} & \color{white}\textbf{Notes}\\\endfirsthead
\rowcolor{CornellReferenceNavy}\color{white}\textbf{Cue / recall prompt} & \color{white}\textbf{Notes (continued)}\\\endhead
\arrayrulecolor{CornellReferenceLineGray}
\textbf{What rule applies to The asm declaration?}\\[-1mm]{\scriptsize\CornellClauseRef{9.10}} & Defines the core-language rules for The asm declaration, including formation, semantic effect, interactions with type and lookup, and any ill-formed or implementation-dependent cases.\\\hline
\end{longtable}
\begin{CornellReferenceSummaryBox}The asm declaration supplies the core-language semantics portion of this clause. The key review move is to connect its local wording to the clause-wide model, then classify each condition by normative force before relying on it.\end{CornellReferenceSummaryBox}
\Needspace{8\baselineskip}
\subsection*{9.11 Linkage specifications}
\begin{longtable}{@{}>{\RaggedRight\arraybackslash\columncolor{CornellReferenceCueGray}}p{0.28\textwidth}>{\RaggedRight\arraybackslash}p{0.64\textwidth}@{}}
\rowcolor{CornellReferenceNavy}\color{white}\textbf{Cue / recall prompt} & \color{white}\textbf{Notes}\\\endfirsthead
\rowcolor{CornellReferenceNavy}\color{white}\textbf{Cue / recall prompt} & \color{white}\textbf{Notes (continued)}\\\endhead
\arrayrulecolor{CornellReferenceLineGray}
\textbf{What rule applies to Linkage specifications?}\\[-1mm]{\scriptsize\CornellClauseRef{9.11}} & Locates names and entity identity for Linkage specifications; distinguish where a name is visible from whether declarations denote the same entity across contexts.\\\hline
\end{longtable}
\begin{CornellReferenceSummaryBox}Linkage specifications supplies the core-language semantics portion of this clause. The key review move is to connect its local wording to the clause-wide model, then classify each condition by normative force before relying on it.\end{CornellReferenceSummaryBox}
\Needspace{8\baselineskip}
\subsection*{9.12 Attributes}
\begin{longtable}{@{}>{\RaggedRight\arraybackslash\columncolor{CornellReferenceCueGray}}p{0.28\textwidth}>{\RaggedRight\arraybackslash}p{0.64\textwidth}@{}}
\rowcolor{CornellReferenceNavy}\color{white}\textbf{Cue / recall prompt} & \color{white}\textbf{Notes}\\\endfirsthead
\rowcolor{CornellReferenceNavy}\color{white}\textbf{Cue / recall prompt} & \color{white}\textbf{Notes (continued)}\\\endhead
\arrayrulecolor{CornellReferenceLineGray}
\textbf{What rule applies to Attribute syntax and semantics?}\\[-1mm]{\scriptsize\CornellClauseRef{9.12.1}} & Defines the recognized forms for Attribute syntax and semantics; syntactic acceptance does not replace the semantic constraints stated with the production.\\\hline
\textbf{What rule applies to Alignment specifier?}\\[-1mm]{\scriptsize\CornellClauseRef{9.12.2}} & Defines the core-language rules for Alignment specifier, including formation, semantic effect, interactions with type and lookup, and any ill-formed or implementation-dependent cases.\\\hline
\textbf{What rule applies to Assumption attribute?}\\[-1mm]{\scriptsize\CornellClauseRef{9.12.3}} & Defines the core-language rules for Assumption attribute, including formation, semantic effect, interactions with type and lookup, and any ill-formed or implementation-dependent cases.\\\hline
\textbf{What rule applies to Carries dependency attribute?}\\[-1mm]{\scriptsize\CornellClauseRef{9.12.4}} & Defines the core-language rules for Carries dependency attribute, including formation, semantic effect, interactions with type and lookup, and any ill-formed or implementation-dependent cases.\\\hline
\textbf{What rule applies to Deprecated attribute?}\\[-1mm]{\scriptsize\CornellClauseRef{9.12.5}} & Defines the core-language rules for Deprecated attribute, including formation, semantic effect, interactions with type and lookup, and any ill-formed or implementation-dependent cases.\\\hline
\textbf{What rule applies to Fallthrough attribute?}\\[-1mm]{\scriptsize\CornellClauseRef{9.12.6}} & Defines the core-language rules for Fallthrough attribute, including formation, semantic effect, interactions with type and lookup, and any ill-formed or implementation-dependent cases.\\\hline
\textbf{What rule applies to Likelihood attributes?}\\[-1mm]{\scriptsize\CornellClauseRef{9.12.7}} & Defines the core-language rules for Likelihood attributes, including formation, semantic effect, interactions with type and lookup, and any ill-formed or implementation-dependent cases.\\\hline
\textbf{What rule applies to Maybe unused attribute?}\\[-1mm]{\scriptsize\CornellClauseRef{9.12.8}} & Defines the core-language rules for Maybe unused attribute, including formation, semantic effect, interactions with type and lookup, and any ill-formed or implementation-dependent cases.\\\hline
\textbf{What rule applies to Nodiscard attribute?}\\[-1mm]{\scriptsize\CornellClauseRef{9.12.9}} & Defines the core-language rules for Nodiscard attribute, including formation, semantic effect, interactions with type and lookup, and any ill-formed or implementation-dependent cases.\\\hline
\textbf{What rule applies to Noreturn attribute?}\\[-1mm]{\scriptsize\CornellClauseRef{9.12.10}} & Defines the core-language rules for Noreturn attribute, including formation, semantic effect, interactions with type and lookup, and any ill-formed or implementation-dependent cases.\\\hline
\textbf{What rule applies to No unique address attribute?}\\[-1mm]{\scriptsize\CornellClauseRef{9.12.11}} & Defines the core-language rules for No unique address attribute, including formation, semantic effect, interactions with type and lookup, and any ill-formed or implementation-dependent cases.\\\hline
\end{longtable}
\begin{CornellReferenceSummaryBox}Attributes supplies the core-language semantics portion of this clause. The key review move is to connect its local wording to the clause-wide model, then classify each condition by normative force before relying on it.\end{CornellReferenceSummaryBox}
\section{Language-Lawyer and Library-Usage Pitfalls}
\begin{CornellReferencePitfallBox}\begin{itemize}[label=\textcolor{CornellReferenceAmber}{$\blacktriangleright$}]
\item Reading complex declarators left-to-right without applying the grammar.
\item Confusing default, value, direct, copy, and list initialization.
\item Assuming a placeholder always deduces references or cv-qualification as intended.
\item Applying an attribute to an entity it cannot appertain to.
\item Assuming a using-directive and using-declaration have the same lookup effect.
\item Treating a note, example, or recommended practice as though it imposed a mandatory program rule.
\end{itemize}\end{CornellReferencePitfallBox}
\section{Programmer and Implementation Checklist}
\begin{CornellReferenceChecklistBox}\begin{itemize}[label=$\square$]
\item Separate specifiers from declarator operators.
\item Determine the declared entity, type, and linkage.
\item Classify the initialization form and narrowing rules.
\item Verify deduction context and placeholder constraints.
\item Check defaulted, deleted, coroutine, and structured-binding special rules.
\item Audit namespace and using effects on lookup.
\item Confirm each attribute's placement and semantic conditions.
\item For \CornellClauseRef{9.1}, verify preamble against the precise rule category and all cited dependencies.
\item For \CornellClauseRef{9.2}, verify specifiers against the precise rule category and all cited dependencies.
\item For \CornellClauseRef{9.3}, verify declarators against the precise rule category and all cited dependencies.
\item For \CornellClauseRef{9.4}, verify initializers against the precise rule category and all cited dependencies.
\end{itemize}\end{CornellReferenceChecklistBox}
\section{Clause Synthesis}
\begin{CornellReferenceOrientationBox}Defines specifiers, declarators, initialization, function definitions, structured bindings, enumerations, namespaces, using declarations, linkage specifications, inline assembly declarations, and standard attributes. Creates and describes entities, their types, initialization, and program-visible properties. A sound review distinguishes syntax and participation from runtime preconditions, keeps notes and examples non-mandatory, and follows cross-clause dependencies before making a portability claim.\end{CornellReferenceOrientationBox}
\section{Self-Test Questions}
\begin{enumerate}
\item What role does Preamble play in the clause's overall model?
\item Which normative distinctions must be kept separate when applying Specifiers?
\item How would you audit a program or implementation that depends on Declarators?
\item Compare Initializers with Function definitions; what different question does each answer?
\item What role does Function definitions play in the clause's overall model?
\item Which normative distinctions must be kept separate when applying Structured binding declarations?
\item How would you audit a program or implementation that depends on Enumerations?
\item Compare Namespaces with The using declaration; what different question does each answer?
\item What role does The using declaration play in the clause's overall model?
\item Which normative distinctions must be kept separate when applying The asm declaration?
\item How would you audit a program or implementation that depends on Linkage specifications?
\item Compare Attributes with General; what different question does each answer?
\item What role does General play in the clause's overall model?
\item Which normative distinctions must be kept separate when applying Storage class specifiers?
\item Scenario: a program relies on an unstated implementation behavior while using Function specifiers. What must be checked before calling the program portable?
\item Scenario: a program relies on an unstated implementation behavior while using The typedef specifier. What must be checked before calling the program portable?
\end{enumerate}
\clearpage\section{Answer Key}
\begin{enumerate}
\item Preamble is one part of the clause's core-language semantics model. Its role is understood by connecting the local rule in 9.1 to the clause orientation and its dependent concepts.
\item Keep formation and well-formedness separate from runtime preconditions and effects. Also distinguish mandatory wording from notes, implementation choice from undefined behavior, and interface declarations from detailed contracts; see 9.2.
\item Audit Declarators by locating the controlling wording in 9.3, listing inputs and type constraints, proving any preconditions, recording effects and failure behavior, and checking lifetime, invalidation, complexity, and synchronization where applicable.
\item Initializers and Function definitions occupy different steps or facility groups. Analyze each under its own subclause, then follow the explicit dependency between them rather than transferring one set of guarantees to the other; begin at 9.4.
\item Function definitions is one part of the clause's operation semantics model. Its role is understood by connecting the local rule in 9.5 to the clause orientation and its dependent concepts.
\item Keep formation and well-formedness separate from runtime preconditions and effects. Also distinguish mandatory wording from notes, implementation choice from undefined behavior, and interface declarations from detailed contracts; see 9.6.
\item Audit Enumerations by locating the controlling wording in 9.7, listing inputs and type constraints, proving any preconditions, recording effects and failure behavior, and checking lifetime, invalidation, complexity, and synchronization where applicable.
\item Namespaces and The using declaration occupy different steps or facility groups. Analyze each under its own subclause, then follow the explicit dependency between them rather than transferring one set of guarantees to the other; begin at 9.8.
\item The using declaration is one part of the clause's core-language semantics model. Its role is understood by connecting the local rule in 9.9 to the clause orientation and its dependent concepts.
\item Keep formation and well-formedness separate from runtime preconditions and effects. Also distinguish mandatory wording from notes, implementation choice from undefined behavior, and interface declarations from detailed contracts; see 9.10.
\item Audit Linkage specifications by locating the controlling wording in 9.11, listing inputs and type constraints, proving any preconditions, recording effects and failure behavior, and checking lifetime, invalidation, complexity, and synchronization where applicable.
\item Attributes and General occupy different steps or facility groups. Analyze each under its own subclause, then follow the explicit dependency between them rather than transferring one set of guarantees to the other; begin at 9.12.
\item General is one part of the clause's core-language semantics model. Its role is understood by connecting the local rule in 9.8.2.1 to the clause orientation and its dependent concepts.
\item Keep formation and well-formedness separate from runtime preconditions and effects. Also distinguish mandatory wording from notes, implementation choice from undefined behavior, and interface declarations from detailed contracts; see 9.2.2.
\item First identify the exact requirement in 9.2.3, then classify the relied-on behavior as required, unspecified, implementation-defined, conditionally supported, or undefined. Portability requires satisfying the program obligations and avoiding dependence on a choice not guaranteed in the target environments.
\item First identify the exact requirement in 9.2.4, then classify the relied-on behavior as required, unspecified, implementation-defined, conditionally supported, or undefined. Portability requires satisfying the program obligations and avoiding dependence on a choice not guaranteed in the target environments.
\end{enumerate}
\section{Key-Term Recap}
\begin{longtable}{@{}>{\RaggedRight\arraybackslash}p{0.28\textwidth}>{\RaggedRight\arraybackslash}p{0.64\textwidth}@{}}\toprule\textbf{Term} & \textbf{Working meaning}\\\midrule\endfirsthead\toprule\textbf{Term} & \textbf{Working meaning (continued)}\\\midrule\endhead\bottomrule\endfoot
declarator & The syntax that combines a name with pointer, reference, array, function, or member structure.\\
initializer & Syntax and semantics that establish an entity's initial state.\\
placeholder type & A type-specifier form whose type is deduced from context.\\
structured binding & A declaration introducing names bound to components of an initializer.\\
linkage specification & A declaration controlling language linkage.\\
attribute & Metadata-like syntax with specified appertainment and semantic rules.\\
\end{longtable}
\CornellTakeaway{A declaration must be analyzed as a grammar-derived type plus entity, scope, linkage, initialization, and attribute effects.}
\end{document}
